\scanpage{25}
\mspage{1}
% archive sheet number: 24
% L1--2.
\uline{Théorie de Rosenlicht}

\textcircled{I} \uline{Existence des Picard pour les courbes}
% L3--9.
Soit $X$ un $Y$-schéma, avec

\begin{enumerate}[label=\alph*)]
\item $Y$ loc. noethérien
\item $X$ propre sur $Y$ ({\unsure{faut-il}} même : loc. projectif ?)
\item $X$ plat sur $Y$
\item Les fibres de $X$ sur $S$ sont des courbes algébriques \uline{absolument} \deleted{{\unsure{connexes}}} \deleted{{\unsure{irréductibles}}} \uline{réduites} et absolument \deleted{réduites}\inserted{connexes} \margin{[ce dernier pt n'est pas essentiel ---]}
\item Le nb de composantes \deleted{\illegible{1}} irréd. absolues des fibres de $X$ sur $Y$ est loc. constant sur $Y$
\end{enumerate}
% slanted
\margin{si on abandonne cette hypothèse, on {\unsure{trouve des|tombe aux}} {\unsure{phénomènes}} à la Néron ?!}
% L10.
Alors le {\unsure{$S$|$f$}}-schéma ${\unsure{\mathscr P_{X/Y}|\mathcal P_{X/Y}}}$ relatif existe, et est \uline{simple} sur $Y$.
% L11--17.
De plus, si le nb des composantes irréductibles absolues est $n$, il y a un revêtement \deleted{$Y$} étale $Y'$ de degré $n$ de $Y$ et si $\mathbb Z_{Y'/Y}$ est le schéma en groupes \inserted{$\prod_{Y'/Y}(\mathbb Z_{Y'})=\underline{\operatorname{Hom}}_Y(Y',\mathbb Z_Y)$} sur $Y$ défini par $Y'$ (qui sur $Y'$ devient isomorphe à $\mathbb Z^n$) on a une augmentation \deleted{{\unsure{surjective}} et {\unsure{propre}}} fidèlement plate
\margin{\uline{{\unsure{références}}}}
\[
\mathscr P_{X/Y}\longrightarrow\mathbb Z_{Y'/Y}\quad\deleted{{\unsure{\infty}}}
\]
dont le noyau sera noté $\mathscr P_{X/Y}^c$, d'où la suite exacte
\[
0\longrightarrow\mathscr P_{X/Y}^c\longrightarrow\mathscr P_{X/Y}\longrightarrow\mathbb Z_{Y'/Y}{\unsure{\longrightarrow|\text{---}}}0
\]
% L18--22.
{\unsure{Noter}} que l'on a un \uline{homomorphisme {\unsure{canon.|canonique}}}
\[
\mathbb Z_{Y'/Y}\longrightarrow\mathbb Z_Y
\]
donc une augmentation
\[
\mathscr P_{X/Y}\longrightarrow\mathbb Z_Y
\]
